% year: תשפ"ה
% type: מבחן
% semester: א
% moed: מיוחד
% number: 3
% points: 16
% categories: taylor
% source: exams/מבחנים/תשפה/מועד ג תשפה.pdf

\begin{question}
חשבו בעזרת קירוב לינארי את $\sqrt[3]{10}$, והעריכו את השגיאה. נמקו.
\end{question}

\begin{hint}
בחרו נקודה קרובה ל-$10$ שבה השורש השלישי ידוע: $x_0 = 8$. העריכו את השגיאה בעזרת שארית לגרנז'.
\end{hint}

\begin{solution}
נגדיר $f(x) = x^{1/3}$ ונבחר $x_0 = 8$ (כי $\sqrt[3]{8} = 2$). הנגזרות:
\[ f'(x) = \tfrac13x^{-2/3}, \qquad f''(x) = -\tfrac29x^{-5/3}, \]
ולכן $f(8) = 2$, $f'(8) = \frac13\cdot\frac14 = \frac1{12}$.

\textbf{קירוב לינארי:} $P_1(x) = 2 + \frac{1}{12}(x - 8)$, ולכן
\[ \sqrt[3]{10} \approx P_1(10) = 2 + \frac{2}{12} = \frac{13}{6} \approx 2.1667. \]

\textbf{הערכת השגיאה:} לפי משפט טיילור עם שארית לגרנז' קיימת $c \in (8, 10)$ כך ש-
\[ \sqrt[3]{10} - \frac{13}{6} = R_1(10) = \frac{f''(c)}{2!}(10 - 8)^2 = -\frac19c^{-5/3}\cdot 4 = -\frac{4}{9}c^{-5/3}. \]
מכיוון ש-$c > 8$ מתקיים $c^{-5/3} < 8^{-5/3} = \frac1{32}$, ולכן
\[ |R_1(10)| < \frac49\cdot\frac{1}{32} = \frac{1}{72} \approx 0.0139, \]
והשגיאה שלילית, כלומר הקירוב גדול מהערך האמיתי.

\textbf{תשובה:} $\sqrt[3]{10} \approx \frac{13}{6} \approx 2.1667$, עם שגיאה קטנה מ-$\frac{1}{72}$; בפרט $\frac{13}{6} - \frac1{72} < \sqrt[3]{10} < \frac{13}{6}$.
לבדיקה: $\sqrt[3]{10} \approx 2.1544$, והשגיאה בפועל כ-$0.0122$.
\end{solution}
